%% contents/introduction.tex
%% Copyright 2021-2022 Tom M. Ragonneau
%
% This work may be distributed and/or modified under the
% conditions of the LaTeX Project Public License, either version 1.3
% of this license or (at your option) any later version.
% The latest version of this license is in
%   http://www.latex-project.org/lppl.txt
% and version 1.3 or later is part of all distributions of LaTeX
% version 2005/12/01 or later.
%
% This work has the LPPL maintenance status `maintained'.
%
% The Current Maintainer of this work is Tom M. Ragonneau.
\chapter{Introduction}
\label{ch:introduction}

\section{Get started}

In the file \verb|thesis.tex|, enter your thesis' title, your name, the date of creation, and the date of revision (if none, put the same as the date of creation).
To compile this document, run in a terminal (for Linux and macOS) the command \verb|make| (necessitates the package latexmk v4.51 or later).

\section{Cross references}

If you have a numbered figure, table, equation, etc, such as
\begin{equation}
    \label{eq:pythagorean}
    a^2 + b^2 = c^2,
\end{equation}
use \verb|\cref{...}| to refer to it.
On this example, it gives~\cref{eq:pythagorean}.
For figures, tables, etc, it will automatically add a name to the reference.

\begin{figure}[ht]
    \centering
    \begin{tikzpicture}
        \draw (0, 0) rectangle (1, 1);
    \end{tikzpicture}
    \caption{A rectangle}
    \label{fig:rectangle}
\end{figure}

For example, to refer to \cref{fig:rectangle}, I only typed \verb|\cref{fig:rectangle}| (without writing explicitely Figure).
If you want several references simultaneously, just separate them with commas in a single \verb|\cref| command.

\section{Glossaries}

In the file \verb|utils/glossary.tex|, put your acronyms and glossaries entries.
The accronyms will be displayed entirely at their first use.
For example, you might need the \gls{kkt} conditions.
This has been typeset as \verb|\gls{kkt}|.
When using in titles, sections, etc, use \verb|\glsfmtshort{...}| and \verb|\glsfmtlong{...}| to have the acronyms and descriptions, respectively.

\section{Index}

I took most of the code from~\cite{higham20}. He describes how to do it carefully.